\documentclass[12pt]{article}

\usepackage{amsmath, amssymb, amsfonts}
\usepackage{geometry}
\usepackage{authblk}
\usepackage{setspace}
\usepackage{cite}
\usepackage{hyperref}

\geometry{margin=1in}
\setstretch{1.2}

\title{\textbf{On the Structural Asymmetry Between General Relativity and Quantum Mechanics:\\
Why Gravity Cannot Be Unified as a Conventional Field}}

\author[1]{Paul William Jarvis}
\affil[1]{Independent Researcher, United Kingdom \\ ORCID: 0009-0009-8933-857X}

\date{2 June 2026}

\begin{document}

\maketitle

\begin{abstract}
General Relativity (GR) and Quantum Mechanics (QM) remain mathematically incompatible despite a century of effort. This paper argues that the failure of unification arises from a structural asymmetry: gravity is not a force within spacetime but the geometric framework that \emph{defines} spacetime. Every particle and field exists within this gravitational structure, making gravity universal in a way no quantum interaction is. Because quantum field theory requires a fixed background while GR denies that such a background exists, unification cannot proceed by treating gravity as another quantum field. This structural asymmetry is named here as the \emph{Gravitational Baseline Principle}, and the paper sets out, in plain terms, how both spacetime and gravity emerge from a deeper substrate that has no geometry of its own.
\end{abstract}

\section{Introduction}

Einstein's General Relativity identifies gravity with the curvature of spacetime itself \cite{einstein1916}. Quantum Mechanics describes matter and radiation through probabilities that unfold on a fixed background of space and time \cite{dirac1930}. Both theories pass every experimental test, yet they do not fit together.

This paper makes one argument: GR and QM cannot be unified as long as gravity is treated as a field analogous to the electromagnetic, weak, or strong interactions. Gravity is not a field \emph{within} spacetime; it is the structure that \emph{defines} spacetime.

\section{Gravity as Universal Geometry}

Einstein's field equations,
\begin{equation}
G_{\mu\nu} = 8\pi G\, T_{\mu\nu},
\end{equation}
express the idea that matter and energy (on the right) determine the curvature of spacetime (on the left). A constant term, written $\Lambda$ and called the cosmological constant, can be added to the left side; it describes the dark energy driving the universe's accelerating expansion. Every particle, regardless of charge or composition, follows the paths this curvature lays out.

By contrast, the quantum interactions are selective:
\begin{itemize}
    \item Electromagnetism acts only on electrically charged particles.
    \item The weak interaction acts only on particles carrying weak charge, and part of it acts only on particles of one handedness.
    \item The strong interaction acts only on quarks and gluons.
\end{itemize}

Gravity alone is universal. It is not one interaction among many; it is the geometric background in which all interactions occur.

\section{The Asymmetry Problem}

Quantum field theory assumes a fixed spacetime background. GR denies that such a background exists. This is a structural contradiction, not a technical one:

\begin{center}
\begin{tabular}{|c|c|}
\hline
\textbf{Quantum Mechanics} & \textbf{General Relativity} \\
\hline
Needs a fixed stage & The stage itself moves \\
Fields play out \emph{on} the stage & The stage \emph{is} the field \\
Effects simply add up & Gravity feeds back on itself \\
\hline
\end{tabular}
\end{center}

Think of matter and fields as hikers crossing terrain. Quantum field theory needs the terrain fixed in advance, so its formulas have somewhere definite to live. General Relativity says the terrain itself sags under whatever is standing on it, and that sag is what then steers every hiker's path; there is no fixed ground to write on in the first place. This incompatibility is central to the failure of every attempt to turn gravity into a quantum theory by treating it as an ordinary field \cite{dewitt1967}.

\section{The Gravitational Baseline Principle}

\textbf{Principle (Gravitational Baseline Principle).}
\emph{Gravity is not a force within spacetime but the geometric structure that defines spacetime. Because all physical systems exist and evolve within this structure, gravity acts as the universal baseline against which all other interactions are measured. Any theory attempting to unify gravity with the quantum forces must therefore work at a deeper level, one that comes before geometry, where spacetime itself is not fundamental.}

This principle names the structural asymmetry between gravity and the quantum interactions: gravity is not one more force alongside the others; it is the underlying geometry that the other forces need in order to be defined at all.

\section{The Substrate Beneath Spacetime}

If gravity is the terrain itself rather than one of the hikers crossing it, unification cannot proceed by adding gravity to the list of quantum forces. Instead, both spacetime and the fields that live in it must come from something prior to geometry altogether.

Picture the universe not as space filled with particles, but as a large collection of elementary quantum systems, correlated with one another to varying degrees. Two systems that are strongly correlated behave as if they are close together; two that are weakly correlated behave as if they are far apart. How strongly two systems are correlated can be measured by how much knowing one tells you about the other, and distance is simply read off from that: more correlation means less distance. Disturbances in the substrate also travel at a fixed, finite speed, which keeps cause before effect as ordinary physics requires, and that speed limit settles what shape the resulting space has. Correlation alone does not: the same pattern of correlation can be read as a flat space or as a saddle-shaped one, depending on how correlation is converted into distance. Only the flat reading gives signals a single, constant speed: read the other way, the same signal would appear to slow down the further it travelled, and there would be no one speed of light at all. The large-scale shape of space then follows the pattern of which systems interact with which, and for a regular pattern it is flat, as the space of our universe is measured to be.

If this is right, spacetime is a large-scale, derived description of an underlying pattern of correlation, in the same way that temperature is a derived description of the average motion of many particles rather than a property carried by any single particle. Gravity would then be the substrate's collective response to the energy it carries, a shared ripple through the whole substrate rather than a change in any one pair of correlations, obeying a large-scale law much as the laws of thermodynamics describe how heat flows without tracking every particle individually.

This is also why gravity resists being turned into a quantum theory in the way the other forces were. Doing that for electromagnetism means writing a field on a background that is already there. Gravity can be treated the same way as a low-energy approximation, treating small ripples in the terrain as a field on a fixed background, and that approximation works. It breaks down at high energies, where it can no longer make predictions, which is exactly what one would expect if the background itself is not fundamental. A successful theory of quantum gravity must start one level deeper, with the substrate itself, rather than starting from spacetime and trying to add quantum behavior on top of it.

\section{Conclusion}

GR and QM resist unification because they describe physics at different conceptual levels. Gravity is not a force within spacetime; it is the structure that defines spacetime. A successful unification must arise from a deeper substrate, one that comes before geometry, from which both spacetime and quantum fields emerge.

\begin{thebibliography}{9}

\bibitem{einstein1916}
A. Einstein, ``The Foundation of the General Theory of Relativity,'' \emph{Annalen der Physik}, 1916.

\bibitem{dirac1930}
P. A. M. Dirac, \emph{The Principles of Quantum Mechanics}, 1930.

\bibitem{dewitt1967}
B. S. DeWitt, ``Quantum Theory of Gravity,'' \emph{Phys. Rev.}, 1967.

\end{thebibliography}

\end{document}
